% Part: first-order-logic
% Chapter: models-theories
% Section: introduction

\documentclass[../../../include/open-logic-section]{subfiles}

\begin{document}

\olfileid{fol}{mat}{int}
\olsection{تعارف}

\begin{explain}
مسلّمی طریقے دی ترقی سائنس دی تاریخ دی اک وڈی کامیابی اے، تے
ریاضی دی تاریخ وچ خاص اہمیت رکھدی اے۔ کسے شعبے نوں مسلّماں
دی بنیاد اُتے قائم کردیاں بہت سارے سوال واضح کرن پیندے نیں:
ایہ شعبہ کس بارے اے؟ ایس دے سبھ توں بنیادی تصور کیہڑے نیں؟
اوہ اک دوجے نال کس طرح جڑے نیں؟ کی شعبے دے سارے تصورات
ایہناں بنیادی تصورات دے ذریعے تعریف کیتے جا سکدے نیں؟ ایہ
تصورات کیہڑے قاعدیاں اُتے عمل کردے نیں، تے کیہڑے قاعدیاں اُتے
لازماً عمل کرنا چاہیدا اے؟

مسلّمی طریقہ تے منطق اک دوجے لئی بہت مناسب نیں۔ رسمی منطق
مسلّمی نظریے بنانے، نظریے دے مسلّماں توں ٹھیک ٹھیک متعین
طریقے نال قضیے ثابت کرن، تے مسلّماں نوں پورا کرن والے سارے
نظاماں دیاں خاصیتاں دی باقاعدہ جانچ لئی وسیلے دیندی اے۔
\end{explain}

\begin{defn}
!!{sentence}s دا سیٹ~$\Gamma$ \emph{بند} اے اُدوں تے صرف اُدوں
جدوں $\Gamma \Entails !A$ ہون اُتے $!A \in \Gamma$ وی ہووے۔
!!{sentence}s دے سیٹ~$\Gamma$ دی \emph{بندش}
$\Setabs{!A}{\Gamma \Entails !A}$ اے۔

اسیں آکھدے آں کہ~$\Gamma$ نوں جملیاں دے سیٹ~$\Delta$ نے
\emph{مسلّماں نال متعین کیتا اے} جے $\Gamma$،~$\Delta$ دی بندش اے۔
\end{defn}

\begin{explain}
اسیں مسلّمی نظریے نوں اوہناں جملیاں دے سیٹ دے طور تے سمجھ سکدے آں
جیہڑا اوہدے مسلّماں دے سیٹ~$\Delta$ توں متعین ہوندا اے۔ دوجے
لفظاں وچ، فرض کرو اسیں اک ایہو جیہی پہلے درجے دی زبان ورتدے
آں جیہدے غیر منطقی نشان اُس سائنس دے بنیادی اجزا لئی نیں
جیہنوں اسیں مسلّمی طریقے نال پڑھنا چاہندے آں، تے نال اوہناں
بنیادی قاعدیاں نوں بیان کرن والے !!{sentence}s دا سیٹ اے۔ تاں
اسیں نظریے نوں ایس زبان دے اوہناں سارے !!{sentence}s نال ظاہر
سمجھ سکدے آں جیہڑے مسلّماں توں منطقی طور تے نکلدے نیں۔
ایہ دائرہ اک ہی بنیادی جزو تے سادہ مسلّماں والے نمونیاں، جیویں
جزوی ترتیباں دے نظریے، توں لے کے نیوٹنی میکانیات ورگے
پیچیدہ نظریاں تک پھیلیا ہویا اے۔

مسلّمی طریقے ول ایس رسمی پہنچ نوں اہم بنان والے منطقی حقائق
ہیٹھاں دتے نیں۔ فرض کرو $\Gamma$ کسے نظریے دا مسلّمی نظام،
یعنی جملیاں دا اک سیٹ، اے۔
\begin{enumerate}
\item اسیں ٹھیک ٹھیک دس سکدے آں کہ مسلّمی نظام کدوں مقصودہ
  !!{structure}s دی جماعت نوں پکڑدا اے۔ یعنی جے ساڈی دلچسپی
  !!{structure}s دی کسے خاص جماعت وچ اے، تاں مسلّمی نظام~$\Gamma$
  اوہ جماعت اُدوں تے صرف اُدوں پکڑے گا جدوں !!{structure}s
  ٹھیک اوہی~$\Struct M$ ہون جیہناں لئی $\Sat{M}{\Gamma}$۔
\item اس کم وچ ناکامی ایس کرکے ہو سکدی اے کہ کجھ $\Struct M$
  ایہو جیہے ہون جیہناں لئی $\Sat{M}{\Gamma}$، پر $\Struct M$
  ساڈے مقصودہ !!{structure}s وچ شامل نہ ہووے۔ ایس توں اسیں
  اوہ مسلّمے شامل کرن ول جا سکدے آں جیہڑے~$\Struct M$ وچ سچے نہ ہون۔
\item جے گھٹ توں گھٹ ایس حد تک کامیابی ہووے کہ $\Gamma$ ساریاں
  مقصودہ !!{structure}s وچ سچا ہووے، تاں جملہ~$!A$ وی ساریاں
  مقصودہ !!{structure}s وچ سچا ہووے گا جدوں وی $\Gamma \Entails !A$۔
  ایس لئی اسیں منطقی وسیلے، جیویں !!{derivation} دے طریقے، ورت کے
  جملیاں دا ساریاں مقصودہ !!{structure}s وچ سچا ہونا صرف ایہ
  دکھا کے ثابت کر سکدے آں کہ اوہ مسلّماں توں منطقی طور تے نکلدے نیں۔
\item کدے ساڈے ذہن وچ مقصودہ !!{structure}s نہیں ہوندیاں، بلکہ
  اسیں خود مسلّماں توں شروع کردے آں: کجھ بنیادی اجزا لیندے
  آں جیہناں لئی اسیں چاہندے آں کہ اوہ کجھ قاعدیاں اُتے عمل کرن، تے انہاں
  قاعدیاں نوں مسلّمی نظام وچ درج کردے آں۔ سبھ توں پہلاں اسیں
  ایہ جانچنا چاہندے آں کہ مسلّمے آپس وچ متضاد نہ ہون؛ جے ہون،
  تاں کوئی تصور انہاں قاعدیاں اُتے عمل نہیں کر سکدا تے ساڈا
  بنایا نظریہ غیر موافق ہووے گا۔~$\Gamma$ دا کوئی نمونہ لبھ کے
  اسیں تصدیق کر سکدے آں کہ ایہ مسئلہ نہیں۔ تے جے ساڈے نظریے
  دے نمونے ہون، تاں اسیں منطقی طریقیاں نال انہاں دی جانچ وی
  کر سکدے آں تے نمونے بنا وی سکدے آں۔
\item مسلّماں دا اک دوجے توں آزاد ہونا وی اہم سوال اے۔ ہو سکدا
  اے کہ کوئی مسلّمہ اصل وچ باقیاں دا منطقی نتیجہ ہووے، تے ایس
  لئی فالتو ہووے۔ اسیں ثابت کر سکدے آں کہ مسلّمہ $!A$،
  جیہڑا $\Gamma$ وچ اے، فالتو اے جے ثابت کر دیئے کہ $\Gamma \setminus
  \{!A\} \Entails !A$۔ ایس دے الٹ، ایہ دکھا کے کہ
  $(\Gamma \setminus \{!A\}) \cup \{\lnot
  !A\}$ پورا ہون جوگا اے، اسیں ثابت کر سکدے آں کہ ایہ مسلّمہ
  فالتو نہیں۔ مثال لئی، ایس طریقے نال دکھایا گیا سی کہ متوازی
  خطاں دا مسلّمہ جیومیٹری دے ہور مسلّماں توں آزاد اے۔
\item اک ہور اہم سوال نظریے وچ تصورات دے تعریف دے قابل ہون
  دا اے: زبان دی چُون ایہ طے کردی اے کہ نظریے دے نمونے کیہناں
  چیزاں توں بندے نیں۔ پر نظریے دے ہر پہلو نوں اوہدے نمونیاں
  وچ الگ توں ظاہر کرنا ضروری نہیں۔ مثال لئی، ہر ترتیب
  $\le$ اک متعلقہ سخت ترتیب~$<$ طے کردی اے: اک مل جائے، تاں
  دوجی تعریف کیتی جا سکدی اے۔ ایس لئی ایہ ضروری نہیں کہ
  ایہو جیہی ترتیب والے نظریے دے نمونے وچ متعلقہ سخت ترتیب
  \emph{وی} الگ شامل ہووے۔ عمومی طور تے اک تعلق دوجیاں دے
  ذریعے کدوں تعریف کیتا جا سکدا اے؟ تے کدوں کسے تعلق نوں
  دوجیاں دے ذریعے تعریف کرنا ناممکن اے، جیہدے نتیجے وچ اوہنوں
  زبان دے بنیادی اجزا وچ شامل کرنا پیندا اے؟
\end{enumerate}
\end{explain}


\end{document}
